\uuid{y0PM}
\exo7id{3358}
\titre{exo7 3358}
\auteur{quercia}
\organisation{exo7}
\datecreate{2010-03-09}
\isIndication{false}
\isCorrection{false}
\chapitre{Matrice}
\sousChapitre{Propriétés élémentaires, généralités}
\cpge{mpsi-mp2i.calcul-matriciel-systemes.e, pcsi.determinants.c, ptsi.determinants.a, mpsi-mp2i.espaces-vectoriels.b, pcsi.espaces-vectoriels.b, ptsi.espaces-vectoriels.a, mpsi-mp2i.dimension-finie.c, pcsi.dimension-finie.c, ptsi.dimension-finie.c, bcpst1.matrices, mp-mpi.structures-algebriques.f, bcpst2.espaces-vectoriels.a}
\module{Algèbre}
\niveau{L1}
\difficulte{}

\contenu{
\question{
Soit $M = (a_{ij}) \in \mathcal{M}_n(K)$. On dit que $M$ est {\it en damier\/} si
$a_{ij} = 0$ pour $j-i$ impair.
On note ${\cal D}$ l'ensemble des matrices $n\times n$ en damier.
Montrer que ${\cal D}$ est une sous-algèbre de $\mathcal{M}_n(K)$.
Quelle est sa dimension ?
}
}
